\documentclass{article}
\usepackage{pathfinder-readable}

\title{Payoff-equivalent feedback in an LLM double auction}

\begin{document}
\maketitle

\begin{abstract}
Q studies language-model traders in a double auction; P studies how peer reports can support
contribution rankings and reward feedback. The thread looked for a way to connect P's methods to the
trades that Q's agents often fail to complete. It found a payoff-equivalent completion signal and two
limits on using P's ranking method with Q's public auction records. The payoff calculation is established
prior work, and no one tested whether the signal changes trading or welfare, so the thread ended at PAUSE.
\end{abstract}

\section{Introduction}

Q, \emph{Competitive Market Behavior of LLMs}, puts 11 language-model buyers and 11 sellers in a double
auction for five rounds \cite{Q}. Each trader knows its own reservation price and sees the standing
offers and order history. A buyer can accept a standing ask by bidding at least that much; a seller can
accept a standing bid in the corresponding way. The agents are told to maximise their trading profit. Q
reports incomplete convergence, especially for GPT Large, whose agents often improve offers by small
amounts until time runs out. Its analysis of Gemini Large reasoning traces links crossing a spread with
words about urgency and execution, but does not establish what causes a trader to cross.

P, \emph{From Truthful Peer Signals to Contribution Rankings}, asks when noisy reports from several
observers can be turned into a common ranking \cite{P}. Its result needs repeated reports about the same
underlying comparison, independent observer errors under stated conditions, and enough overlap between
observers. It also discusses potential-based reward shaping: intermediate feedback can change while the
total reward difference between complete paths stays fixed. P's own labels ``Q'' and ``P'' refer to other
papers on mechanism design and cooperative credit, rather than to the assigned auction paper and P
itself.

The thread first asked whether P's peer calibration could diagnose Q's market. That route lacked the
reports P requires. The peers then pursued a narrower question: could a temporary signal for completing a
trade change the behaviour of Q's fixed language-model traders while leaving final monetary profit
unchanged?

\section{What was done}

The peers checked the identity of the assigned papers before trying to combine their claims. They
compared P's requirement for repeated, conditionally independent labels of one realised comparison with
Q's strategic bids, asks and shared public history. They found no direct input to P's calibration method
in the auction record. They also constructed two private-value assignments that give the same public
trades but different realised welfare.

They next took P's reward-shaping calculation and adapted it to Q's undiscounted trading profit. They
proposed provisional feedback when a trader completes a trade, followed by a full reversal at settlement.
They checked the resulting payoff path by path and compared the proposed claim with prior work on
multi-agent reward shaping \cite{devlin2011,lu2011}. Their source search also found a related LLM-guided
double-auction shaping study \cite{related2025}; it did not establish novelty for this proposed
intervention.

Finally, they tightened the proposed test after two objections. Q's prompts do not explicitly tell agents
the numerical order cap or identify a final turn. Also, actions in an earlier round can affect later
public histories. The peers therefore specified a disclosed final opportunity in the fifth round for a
local decision test, and separate paired full-market runs for a welfare test. They checked that Q's
welfare measure is realised surplus, which need not move with trade count or price dispersion.

\section{Results}

The payoff calculation is exact for a completed round. Let \(i\) denote a trader, \(t\) a step in the
round, and \(T\) its settlement step. Let \(U_i\) be the trader's total profit without added feedback.
Choose a feedback balance \(\psi_{i,t}\) that starts and ends at zero on every possible path. Give the
trader the change \(F_{i,t}=\psi_{i,t+1}-\psi_{i,t}\) at each step. Then
\[
\widetilde U_i
=U_i+\sum_{t=0}^{T-1}F_{i,t}
=U_i+\psi_{i,T}-\psi_{i,0}
=U_i.
\]
Here \(\widetilde U_i\) is final profit including the feedback. A temporary balance can rise when a trade is completed and return to zero at
settlement. Because every complete path has the same final profit as before, profit best responses and
Nash equilibria stay the same. This is P's telescoping argument in Q's setting, not a new equilibrium
result \cite{P,devlin2011,lu2011}. A change in fixed agents' choices under the displayed feedback would
be a behavioural finding about the feedback, not a change in the game's monetary incentives.

The peers' welfare check also holds for Q's stated values. The five buyer--seller pairs with strictly
positive gains can produce surplus \(2.50+2.00+1.50+1.00+0.50=7.50\), the full available surplus. The
sixth pair has value and cost \(2.00\), so it adds no surplus. Thus five trades can achieve full
allocative efficiency even if their prices differ. A test of market improvement must measure realised
surplus; crossing rates, trade count and price dispersion can help explain it but do not determine it
\cite{Q}.

The public record alone cannot always reveal that surplus. In the peers' example, a seller with cost
\(0.75\) trades with buyer A at price \(2.00\), while buyer B does not trade. A has value \(3.25\) and B
has \(2.25\) in one assignment; the values are swapped in another. The same public actions are feasible
in both, and the best attainable surplus is unchanged, yet this trade creates \(2.50\) surplus in the
first assignment and \(1.50\) in the second. Q's experimenter can compute welfare because it knows the
private values. The example limits a diagnostic based only on the public tape; it does not challenge Q's
reported efficiency calculation.

\section{Objections and limits}

P's ranking result needs reports from observers of the same underlying comparison, with errors that are
independent once that comparison is fixed. Q records traders' strategic actions under a shared order
history, not those reports. The peers illustrated the problem with a hypothetical public anchor that
every observer repeats regardless of the true comparison: perfect agreement would then convey no truth.
This is an objection to applying P's assumptions to the auction tape, not an empirical finding that P's
method fails on Q's data.

Nor does Q's late small-step quoting prove that its agents passed up a known final profitable trade. The
displayed prompts give the round number and mention an order cap, but not the numerical cap or a
countdown. For a strict local comparison, the peers required the final actor in the fifth and final round
to be told that no later action can fill an unaccepted quote. If a resting opposite-side offer yields
positive profit when crossed, crossing then earns more than an unfilled offer. Earlier rounds do not give
this simple comparison because today's actions may alter tomorrow's history.

No completion-feedback intervention was run. The algebra therefore gives no estimate of its effect on
crossing or welfare. The peers' source check also placed payoff invariance in prior work
\cite{devlin2011,lu2011}. The related auction study found in the search was noted, but the available
material did not establish whether it tests this same matched feedback question \cite{related2025}.

\section{Why the thread stopped}

The verifier paused the thread because the supported payoff identity was already known and the proposed
behavioural effect had not been measured. The peer-calibration and public-tape checks set useful
boundaries, but did not supply an intervention result for the assigned pair. The smallest restart is a
repeated, matched final-action probe: disclose the last opportunity and the same profitable standing
offer in both conditions, keep the state and final profit fixed, and compare crossing with and without
provisional completion feedback and its stated settlement reversal. A claim about market welfare would
then need paired full-market runs with realised surplus as the main outcome.

\bibliographystyle{plainurl}
\bibliography{references}

\end{document}
